\documentclass{article}
\usepackage{amsmath, amssymb, amsthm}

\setlength{\parindent}{0pt}

\newtheorem{claim}{Claim}

\begin{document}

\begin{claim}
    There are $15120$ ways to arrange BANANARAMA.
\end{claim}

\begin{proof}
    The word BANANARAMA consists of $10$ letters, with $5$ A's, $2$ N's, and one each of B, R, and M.
    Since identical letters are indistinguishable, the number of distinct arrangements is
    \[
    \frac{10!}{2! \cdot 5!} = 15120
    \]

    \medskip

    Hence there are $15120$ ways to arrange BANANARAMA.
\end{proof}

\begin{claim}
    There are $\frac{72!}{12! \cdot 24! \cdot 36!}$ different paths from point $(0, 0, 0)$ to 
    $(12, 24, 36)$, assuming every step affects only one coordinate.
\end{claim}

\begin{proof}
    The entire path consists of $72$ steps, with $12$ steps in the $x$ direction, $24$ steps in the $y$ 
    direction, and $36$ steps in the $z$ direction. Since steps in the same coordinate are 
    indistinguishable, the number of distinct paths is
    \[
    \frac{72!}{12! \cdot 24! \cdot 36!}
    \]

    \medskip

    Hence there are $\frac{72!}{12! \cdot 24! \cdot 36!}$ different paths from point $(0, 0, 0)$ to 
    $(12, 24, 36)$, assuming every step affects only one coordinate.
\end{proof}

\begin{claim}
    There are $\frac{(2n)!}{2^n \cdot n!}$ different ways to pair up $2n$ students.
\end{claim}

\begin{proof}
    There are $(2n)!$ linear orderings of $2n$ students. Reading each ordering as consecutive pairs 
    produces a pairing. However, this procedure overcounts each pairing. Firstly, the two students 
    within each pair can be interchanged, giving a factor of $2^n$. Secondly, the $n$ pairs 
    themselves can be permuted arbitrarily, giving an additional factor of $n!$. Hence each pairing 
    is overcounted by $2^n \cdot n!$, so the number of distinct pairings is
    \[
    \frac{(2n)!}{2^n \cdot n!}
    \]

    \medskip

    Hence there are $\frac{(2n)!}{2^n \cdot n!}$ different ways to pair up $2n$ students.
\end{proof}

\begin{claim}
    
\end{claim}

\begin{proof}
    
\end{proof}

\end{document}
